\documentclass{article}

\usepackage{booktabs}
\usepackage{tabularx}

\title{Development Plan\\\progname}

\author{\authname}

\date{}

\input{../Comments}
%% Common Parts

\newcommand{\progname}{Software Engineering} % PUT YOUR PROGRAM NAME HERE
\newcommand{\authname}{Team 8, AlgoCatan
\\ Jake Read
\\ Rebecca Di Filippo
\\ Matthew Cheung
\\ Sunny Yao} % AUTHOR NAMES

\usepackage{hyperref}
    \hypersetup{colorlinks=true, linkcolor=blue, citecolor=blue, filecolor=blue,
                urlcolor=blue, unicode=false}
    \urlstyle{same}
                                


\begin{document}

\maketitle

\begin{table}[hp]
\caption{Revision History} \label{TblRevisionHistory}
\begin{tabularx}{\textwidth}{llX}
\toprule
\textbf{Date} & \textbf{Developer(s)} & \textbf{Change}\\
\midrule
Sept. 22 & Sunny,Rebecca & inital dev plan created for milestone 1\\
April 6 & Rebecca & Final revision\\
\bottomrule
\end{tabularx}
\end{table}

\newpage{}

\raggedright
This document outlines the development plan for RLCatan, an
AI agent that learns to play the board game \emph{Catan} competitively 
using reinforcement learning. 

\section{Confidential Information}

\raggedright
There is no confidential information to protect in this project.

\section{IP to Protect}

\raggedright
There is no IP to protect in this project.

\section{Copyright License}

\raggedright{AlgoCatan is adopting the MIT License, which can be found at \href{https://github.com/AlgoCatan/RLCatan/blob/main/LICENSE}{this link}.}



\section{Team Meeting Plan}

To maintain project momentum and ensure high-quality communication, the team adheres to a structured synchronization protocol consisting of weekly status updates and targeted technical sessions.

\begin{itemize}
    \item \textbf{Standing Synchronous Meetings:} Held every Thursday from 4:30 PM -- 6:30 PM. These sessions are dedicated to high-level project tracking, addressing technical blockers, and redistributing tasks based on the current sprint's progress.
    \item \textbf{Meeting Structure \& Documentation:} 
    \begin{itemize}
        \item \textbf{Agenda Management:} The designated team organizer distributes a formal agenda via Discord 24 hours prior to the meeting. This includes "Action Items" carried over from the previous week and "New Items" for immediate discussion.
    \end{itemize}
    \item \textbf{Communication Channels:} 
    \begin{itemize}
        \item \textbf{Discord:} Serves as the primary hub for daily asynchronous coordination.
        \item \textbf{Physical Presence:} In-person meetings are held in the Engineering Technology Building (ETB) to facilitate collaborative whiteboarding, with hybrid support available via video link.
    \end{itemize}
    \item \textbf{Supervisor Consultations:} Formal meetings with the Project Supervisor are scheduled every Monday at 5:15-6:00 PM, online through Zoom. These sessions focus on high-level project guidance, milestone reviews, and strategic feedback.
\end{itemize}

\section{Team Communication Plan}


We will be using Discord as our main communication 
platform. We will create a server with channels for general 
discussion, meeting scheduling, and project management.
For more formal communication, we will use email to contact 
our advisor.
We will also use GitHub Issues for tracking existing 
issues and their closures.

\newpage{}
\section{Team Member Roles}

To ensure efficient development, team members are assigned specific focus areas. While everyone will collaborate, these roles define primary ownership of the project’s core components.

\begin{itemize}
  \item \textbf{Jake (Project Lead + Integration Engineer + AI design):} 
    \begin{itemize}
        \item \textbf{Technical:} Core AI agent design and Python backend development.
        \item \textbf{Non-Technical:} Project management and primary contact for the TA/Supervisor.
    \end{itemize}
    
  \item \textbf{Becca (Organizer + AI optimization + Full-stack):} 
    \begin{itemize}
        \item \textbf{Technical:} Agent optimization and User interface Development.
        \item \textbf{Non-Technical:} Managing agendas for meetings and to do lists. Keeps team up to date with deadlines and deliverables.
    \end{itemize}

  \item \textbf{Sunny (IT Specialist + Devops + Full-stack):} 
    \begin{itemize}
        \item \textbf{Technical:} DevOps (CI/CD, Docker) and UI component implementation. General backend engineering support.
        \item \textbf{Non-Technical:} GitHub repository management and environment troubleshooting, ensuring all team members have access to necessary tools and resources.
    \end{itemize}

  \item \textbf{Matthew (Researcher + Backend + Architecture):} 
    \begin{itemize}
        \item \textbf{Technical:} System architecture and general backend development.
        \item \textbf{Non-Technical:} Researching AI strategy papers and competitive data.
    \end{itemize}
\end{itemize}

As the project progresses, these roles may shift to ensure a balanced workload.


\section{Workflow Plan}


The project follows a professional \textbf{Feature Branch Workflow} to ensure the \texttt{main} branch remains stable, deployable, and verified at all times.

\subsection{GitHub Workflow and Branching Strategy}
\begin{itemize}
    \item \textbf{Branch Naming Conventions:} All development must occur on dedicated branches following the format:
    \begin{itemize}
      \item \textbf{issue\#-short-description} (e.g., \texttt{issue\#42-add-reward-shaping})
    \end{itemize}
    \item \textbf{Pull Requests (PRs) \& Peer Review:} No code is merged directly into \texttt{main}. Every PR requires at least one peer review. Reviewers evaluate logic, adherence to the project's style guides (PEP 8 for Python), and the inclusion of necessary unit tests.
    \item \textbf{Conflict Management:} To maintain a clean, readable project history, developers are encouraged to rebase their feature branches against \texttt{main} prior to final merging.
\end{itemize}

\subsection{Issue Tracking and Project Management}
\begin{itemize}
    \item \textbf{Standardized Reporting:} We utilize GitHub Issue Templates for both "Bug Reports" and "Feature Requests." This ensures that every task has a clear way of describing the problem, expected behavior, and steps to reproduce (for bugs) or a detailed description of the desired functionality (for features).
    \item \textbf{Labeling Architecture:} Issues are categorized using a multi-tag system covering scope (\texttt{backend}, \texttt{frontend}, \texttt{RL-env}), priority (\texttt{high}, \texttt{blocker}), and status (\texttt{in-progress}, \texttt{review-required}).
\end{itemize}

\subsection{Continuous Integration and Quality Assurance (CI/QA)}
\begin{itemize}
    \item \textbf{Automated Testing (CI):} GitHub Actions are configured to trigger our test suite on every PR. This includes \texttt{pytest} for backend logic and \texttt{Vitest} for frontend state management. Merging is restricted if the CI pipeline fails.
    \item \textbf{Static Analysis:} \textbf{SonarQube} and \textbf{ESLint} are integrated into the pipeline to monitor code smells, technical debt, and potential security vulnerabilities.
    \item \textbf{Deployment (CD):} The project is containerized using \textbf{Docker}. Our deployment pipeline automates the delivery of stable builds to our dedicated computing node, ensuring the training environment for our RL agent remains consistent with the development environment.
\end{itemize}

\section{Project Decomposition and Scheduling}


\raggedright
We will be using GitHub Projects to manage our project tasks and 
track milestones. Our project Kanban board will have columns for ``To Do'', ``In Progress'', ``In Review'',
and ``Done'', with each task represented as a card that moves across the columns as it progresses. Tasks
and milestones will be scheduled according to the deadlines outlined in the course, ensuring that
the team meets all deliverable requirements on time. The teams Notetaker is resonsible for keeping
this board up to date. Large tasks will be broken down into smaller,
manageable subtasks to make progress easier to track and estimate. Task assignments will be based
on each team member’s strengths, experience, and availability. To use our resources realistically,
we will estimate the time required for each task, monitor progress. If needed we will adjust the
schedule or reassign tasks.

\medskip
The link to our GitHub project board can be found \href{https://github.com/AlgoCatan/RLCatan}{here}.

\medskip
The major milestones for our project are as follows:
\begin{table}[h!]
  \centering
  \begin{tabular}{|l|l|}
  \hline
  \textbf{Milestone} & \textbf{Due Date} \\ \hline
  Team Formed, Project Selected & Sept. 15 \\ \hline
  Problem Statement, POC Plan, Development Plan & Sept. 22 \\ \hline
  Requirements Document \& Hazard Analysis Revision & Oct. 6 \\ \hline
  V\&V Plan Revision & Oct. 27 \\ \hline
  Design Document Revision & Nov. 10 \\ \hline
  Proof of Concept Demonstration & Nov. 17 \\ \hline
  Design Document Revision & Jan. 19 \\ \hline
  Revision 0 Demonstration & Feb. 2 \\ \hline
  V\&V Report \& Extras Revision 0 & March 9 \\ \hline
  Final Demonstration (Revision 1) & March 23 \\ \hline
  EXPO Demonstration & TBA \\ \hline
  Final Documentation (Revision 1) & April 6 \\ \hline
  
  \end{tabular}
  \caption{Project Milestones and Due Dates}
  \label{tab:project-milestones}
  \end{table}


\section{Proof of Concept Demonstration Plan}

\raggedright
The main risks for the success of our project are the AI not 
being able to effectively learn to play \emph{Catan} at a high level, 
and the user interface not being intuitive or user-friendly. 

To mitigate these risks, we will conduct benchmark testing
after each training run against existing algorithmic bots.
This will allow us to track the AI's performance over time and
make adjustments to the training process as needed.
We will also conduct parallel training runs with different
hyperparameters to identify the most effective training 
strategies.

When we demonstrate our proof of concept,
we will showcase that the AI is able to read gamestate and return valid moves
in \emph{Catan}, even if it is not performing at a high level at this stage. 

\section{Expected Technology}

The selection of the \progname{} technology stack is based on a trade-off analysis between development velocity, library ecosystem support, and performance requirements for real-time reinforcement learning and computer vision.

\subsection{Programming Languages and Core Libraries}
\textbf{Python} was selected as the primary language for the backend and AI development. 

\begin{itemize}
    \item \textbf{Pros:} Extensive ecosystem for Reinforcement Learning RL and Computer Vision CV; seamless integration with \texttt{PyTorch} and \texttt{OpenAI Gym}.
    \item \textbf{Cons:} Slower execution speed compared to compiled languages like C++.
    \item \textbf{Comparison:} While C++ offers superior performance for game simulations, Python’s development speed and the availability of the \textbf{Catanatron} library, which provides a high-performance, Gym-compatible interface for \textit{Catan}, outweighing the raw execution benefits of a lower-level language for this project's scope.
\end{itemize}

\subsection{Frontend and Digital Twin}
We will utilize \textbf{React} and \textbf{TypeScript} for the web-based User Interface.

\begin{itemize}
    \item \textbf{Pros:} Component-based architecture is ideal for building a DigitalTwin that requires real-time state synchronization; strong support for \texttt{WebSocket} integration.
    \item \textbf{Cons:} Higher initial setup complexity compared to vanilla JavaScript or lightweight libraries like Alpine.js.
    \item \textbf{Comparison:} Unlike Angular or Vue, React’s virtual DOM and hooks pattern allow for more efficient re-rendering of complex board states. This is critical for providing a smooth experience during fast-paced AI move suggestions and real-time state updates.
\end{itemize}

\subsection{Computer Vision Stac (DEFFERED)}
For image processing and object detection, the project will implement \textbf{OpenCV} and \textbf{YOLOv9}.

\begin{table}[h!]
\centering
\caption{Comparison of Computer Vision Technologies}
\label{tab:cv-tech-comparison}
\begin{tabularx}{\textwidth}{|l|X|X|}
\hline
\textbf{Technology} & \textbf{Pros} & \textbf{Comparison / Alternative} \\ \hline
\textbf{YOLOv9} & State-of-the-art accuracy and real-time inference speed. & Superior to YOLOv8 and Faster R-CNN due to Programmable Gradient Information (PGI), which improves detection of small board assets. \\ \hline
\textbf{OpenCV} & Industry standard for image pre-processing and calibration. & Preferred over PIL or Scikit-Image due to robust support for camera calibration and perspective transforms required for skewed board views. \\ \hline
\end{tabularx}
\end{table}

\subsection{DevOps and Automation}
We will implement \textbf{GitHub Actions} for our CI/CD pipelines.

\begin{itemize}
    \item \textbf{Pros:} Deep integration with our existing repository; zero-cost infrastructure for student organizations; simple YAML-based configuration.
    \item \textbf{Cons:} Limited customization for highly complex deployment environments compared to self-hosted solutions.
    \item \textbf{Comparison:} While Jenkins offers more plugin-based customization, GitHub Actions eliminates the need for managing a separate build server. This allow the team to focus entirely on software development rather than infrastructure maintenance, which is vital given the eight-month project timeline.
\end{itemize}

\section{Coding Standard}

Since our projects backend is primarily in Python, we will be following the PEP 8 Coding
Standard. The link to this standard can be found
\href{https://peps.python.org/pep-0008/}{here}.

\medskip

Our frontend is primarily in JavaScript and React, so we will be following the google style guide
coding standard. The link to this standard can be found \href{https://google.github.io/styleguide/jsguide.html}{here}.

\newpage{}

\section*{Appendix --- Reflection}

%\input{../Reflection.tex}

\begin{enumerate}
    \item Why is it important to create a development plan prior to starting the
    project?
    \item In your opinion, what are the advantages and disadvantages of using
    CI/CD?
    \item What disagreements did your group have in this deliverable, if any,
    and how did you resolve them?
\end{enumerate}

\subsection*{Jake Read:}\label{subsec:jake-read-reflection}
\begin{enumerate}
    \item
        A development plan helps formulate our expectations for the work in the coming months.
        It allows us to pre-plan our strategies for various milestones and encourages us to look ahead into what technology we’ll be using at each design stage.
        I feel it goes hand in hand with the problem statement in regard to preventing scope creep, as it details an initial outline we planned to follow from initiation of the project.
        Early decisions regarding elements such as the coding standard we’ll use will ensure that our design philosophy remains consistent throughout the project and prevent conflicting styles.
        Additionally, establishing confidential info, IPs to protect, copyright licenses, etc. as early as possible is critical for making decisions going forward.

    \item
        The main advantage of CI/CD in my opinion is automated testing.
        Having the full test suite run on new builds can help catch bugs much earlier and leads to improved code quality.
        It’s also faster to release new builds once all the pipelines are set up, since the process is fully automated, and this can also improve deployment consistency since there’s less human error.
        As for the disadvantages, while some of us have used CI/CD in co-ops, none of us have actually built any CI/CD pipelines before, so it will take significant effort to set it all up, and we’re likely to make mistakes that will cause deployment issues.
        It’s still probably worth it though, since using it is the only way we’ll get that experience.

    \item
        We did not come to any disagreements during this deliverable.
\end{enumerate}

\subsection*{Rebecca Di Filippo:}\label{subsec:rebecca-difilippo-reflection}
\begin{enumerate}
  \item Creating a development plan before starting a project helps the team organize the scope,
   timeline, and responsibilities. It acts as a roadmap so we don’t waste time
    moving in the wrong direction. The developmet plan also makes sure everyone is on the same page 
    about attendance, meetings, roles and responsibilities. Without a plan, there is possibility of miscommunication,
    disorganization, and going in the wrong direction.
  \item The biggest advantage of CI is that it automatically runs tests when you make a commit
   or open a pull request. That way, you know right when you commit if the code is safe to merge, which helps 
   avoid integration problems and conflicts. CD makes deploying new versions faster and more reliable,
    which saves time and helps avoid mistakes. The downside to CI/CD is that it can take some effort,
    and it’s annoying when builds are slow.
  \item Our group didn’t have major disagreements. The closest thing to a disagreement was, at first we
  thought CI/CD would be too difficult and unnecessary for the project. However, after discussing it
  with Dr. Istvan David, we concluded that it is worth including given the scope of our project. Its also
  a good learning experience.
\end{enumerate}

\subsection*{Sunny Yao:}\label{subsec:sunny-yao-reflection}
\begin{enumerate}
  \item I have the opinion that initial planning is always highly important.
  It gets the team on the same page and sets expectations for the project.
  Most importantly, it helps set the vision for the project and define project scope.
  This lets members think about the technologies that will be used and the potential challenges.
  \item An advantage of CI/CD is that it automates initial testing of future
commits and pull requests. This helps catch bugs early and avoid builds crashing in production.
The main disadvantage is the maintenance overhead from the setup as well as the
time barrier needed for each pull request. This can then incentivize
developers to make large infrequent commits to reduce the number of times that
the CI/CD pipeline needs to be run.
  \item We disagreed on initial project selection but each member was given
a chance to make their case through votes and meeting discussions. This helped to
reach a consensus on project direction. In addition, we were able to meet with our
supervisor to get further motivations to pursue our final project idea.
\end{enumerate}

\subsection*{Matthew Cheung:}\label{subsec:matthew-cheung-reflection}
\begin{enumerate}
  \item Making a development plan is important because it acts as a roadmap for the entire project.
  It makes the team to sit down and explicitly think about the next steps we want to take to reach our goal.
  Without it, we would start coding without direction and end up wasting a ton of time developing the wrong things.
  The plan also helps us limit the scope of the project, setting distinct boundaries and limits to whatever work we do.

  \item The biggest advantage for CI/CD is how it automates tedious tasks in the project.
  By automatically building and testing our code it saves the team valuable time during development.
  This catches bugs earlier and ensures of no integration problems when we try to merge code, keeping the project stable.
  However, the downside is the overhead costs it takes to set up.
  It can be a lot of work to get the pipelines running, and for a small team like ours, the time could be costly.
  Additionally, if the builds are slow, it can become annoying and slow down our workflow.

  \item We had some initial disagreements for our project choice, where we were debating between two different projects, but we ended up being force to choose Catan anyway due to prof availability so it didn't really matter.
  Once we got over the initial hurdle and started breaking down the work, we all got on the same page and were ready to work.

\end{enumerate}

\section*{Appendix --- Team Charter}


\subsection*{External Goals}

%\wss{What are your team's external goals for this project? These are not the
%goals related to the functionality or quality of the project.  These are the
%goals on what the team wishes to achieve with the project.  Potential goals are
%to win a prize at the Capstone EXPO, or to have something to talk about in
%interviews, or to get an A+, etc.}

\raggedright
Our teams external goals for this project are to attain a strong understanding
of reinforcement learning and computer vision, to be able to build a hireable portfolio. 
We also are aiming to get a A/A+ in the course, since we plan to put a lot of effort into the
project.

\subsection*{Attendance}

\subsubsection*{Expectations}

%\wss{What are your team's expectations regarding meeting attendance (being on
%time, leaving early, missing meetings, etc.)?}

\raggedright
Our team meets every week at 4:30-6:30 pm on Thursday. All other meetings are 
scheduled on a weekly basis, depending on deadlines and amount of work. It is expected that all
team members attend every meeting and arrive on time. If a team member is unable
to attend a meeting they must notify the team at least 12 hours in advance (withholding an emergency situation). If a 
team member is going to be late or must leave early they must notify the team prior
to the meeting. 

\subsubsection*{Acceptable Excuse}

%\wss{What constitutes an acceptable excuse for missing a meeting or a deadline?
%What types of excuses will not be considered acceptable?}

\raggedright
An acceptable excuse for missing a meeting or deadline includes sickness, family
emergencies, or academic obligations. On the other hand, unacceptable excuses include
forgetting, being lazy, or not prioritizing the team. Again, it is expected that the team
knows 12 hours in advance if a member is going to miss a meeting or deadline. This time
frame allows the team to adjust their plans accordingly.

\subsubsection*{In Case of Emergency}

%\wss{What process will team members follow if they have an emergency and cannot
%attend a team meeting or complete their individual work promised for a team
%deliverable?}

\raggedright
In case of an emergency, the team member must notify the team as soon as possible, before
the meeting or deadline. If a deadline cannot be met, other members will distribute 
the work among themselves to ensure the deadline is met. Its important to note that 
1-2 missed deadline with little notice is acceptable, but repeated offenses will not 
be tolerated.

\subsection*{Accountability and Teamwork}

\subsubsection*{Quality} 

%\wss{What are your team's expectations regarding the quality
%of team members' preparation for team meetings and the quality of the
%deliverables that members bring to the team?}

\raggedright
It is expected that all team members come to meetings prepared with the assigned
task completed for each week. It is expected that all tasks we complete are of high quality
and have been reviewed by at least one other team member before submission. Each
deliverable will be reviewed by all team members before submission to ensure quality 
and consistency. Lastly, all coding should follow the agreed upon coding standards.


\subsubsection*{Attitude}

%\wss{What are your team's expectations regarding team members' ideas,
%interactions with the team, cooperation, attitudes, and anything else regarding
%team member contributions?  Do you want to introduce a code of conduct?  Do you
%want a conflict resolution plan?  Can adopt existing codes of conduct.}

\raggedright
The team leader is responsible for ensuring that all team members are treated fairly.
All team members need to be open to new ideas and willing to compromise. If conflicts
arise, the team will discuss the issue and come to a resolution that works for everyone.
If a team member is not contributing to the team, the team leader will address the issue
with the member privately. If the issue persists, the team will discuss the issue with
the TA or instructor.

\subsubsection*{Stay on Track}

Our team will stay on track by setting clear weekly goals, 
using GitHub Kanban boards to monitor progress, and
holding regular check-ins to ensure accountability. Tasks 
will be distributed fairly according to skills and proficiencies, 
and progress will be tracked through version control commits and 
documented updates. Members who perform well will be recognized 
for their contributions and can take greater lead on project
direction in future tasks. If a member's performance 
falls below expectations, we will first address the issue 
through direct communication and support, ensuring that 
individual effort is accurately reflected in evaluations.
We expect to have high attendance from all members with prior 
warning for any absences.

\subsubsection*{Team Building}

%\wss{How will you build team cohesion (fun time, group rituals, etc.)? }

\raggedright
Considering our team will be working together for 8 months, its important to build
team cohesion. We will do this by scheduling fun activities outside of meetings. one
of those fun activities will be a team dinner at the end of the project to celebrate.
We also plan to play \textit{Catan} together outside of meetings. 

\subsubsection*{Decision Making} 

%\wss{How will you make decisions in your group? Consensus?  Vote? How will you
%handle disagreements? }

\raggedright
Our team will converse together before making any major decisions. We will try to reach
a consensus, but if we are unable to do so we will vote through discord polls. In the event
there is a disagreement , the team leader will have the final say.

\end{document}